\chapter*{JUSTIFICACIÓN DETALLADA E HIPÓTESIS}
\addcontentsline{toc}{chapter}{JUSTIFICACIÓN DETALLADA E HIPÓTESIS}
\chaptermark{JUSTIFICACIÓN E HIPÓTESIS}

\section*{1. Justificación e Importancia}

\textbf{Justificación Práctica y Tecnológica:} Proporciona a los equipos de
desarrollo y arquitectos de software un marco de control cuantitativo para
visibilizar y medir el costo real de mantenimiento, la complejidad oculta y
la vulnerabilidad del código generado por IAG en aplicaciones del sector
asegurador.

\textbf{Justificación Académica y Científica:} Cubre un vacío en la
literatura mediante la formulación matemática del Índice de Deuda Técnica
Generativa (IDTG) y la categorización de una nueva taxonomía de antipatrones
producidos por LLMs.

\section*{2. Formulación de Hipótesis}

\textbf{Hipótesis General ($H_1$):} El desarrollo de aplicaciones
empresariales en Seguros y Reaseguros con uso intensivo de IAG incrementa de
forma estadísticamente significativa la Deuda Técnica Acumulada
($TDR \ge 35\%$) y reduce el Maintainability Index ($MI \le 20$), en
comparación con el desarrollo tradicional.

\textbf{Hipótesis Específica 1 ($H_{1.1}$):} El código en módulos de
tarificación desarrollado con IAG presenta una Complejidad Ciclomática y
Cognitiva $v(G) > 10$ en más del 40\% de sus funciones.

\textbf{Hipótesis Específica 2 ($H_{1.2}$):} Los componentes de Reaseguros
presentan un incremento de al menos 25\% en duplicación de código y un alto
acoplamiento ($CBO \ge 15$).

\textbf{Hipótesis Específica 3 ($H_{1.3}$):} El tiempo de resolución de
cambios ($TTR\text{-}CR$) y corrección de errores ($TTFB$) en reglas
aseguradoras desarrolladas con IAG es un 30\% mayor debido a la baja cohesión
y legibilidad.
